\section{Experimental Results}
\label{sec:ExpResults}

\subsection{Experimental Setup}
\label{sec:ExpSetup}

% Wolverine
% AMD Xilinx Alveo U280
% Intel(R) Xeon(R) CPU E5-2640 v3 @ 2.60GHz

% CPU Experiments
% Intel(R) Core(TM) i9-9900 CPU @ 3.10GHz


% OCT
% AMD Xilinx Alveo U280
% Intel(R) Xeon(R) Gold 6226R CPU @ 2.90GHz
% 100 GbE data center switches
% CMAC/ UltraScale+ Integrated 100G Ethernet Subsystem
% UDP/ Xilinx 
% Ethernet Controller XL710 for 40GbE QSFP+
We compare our FPGA-based MiniRocket, MultiRocket, and Hydra inference engines to software implementations written in (multi-threaded) Python 3.10 in the Aeon-toolkit \cite{aeon-arxiv, aeon-pkdd}. All models were trained using Aeon, and all experiments were performed using the three largest time series from the UCR Time Series Archive \cite{UCRTimeSeries}: \textit{InsectSound} (Training size: 139883, Test set size: 139883, Length: 3750), \textit{MosquitoSound} (Training size: 25000, Test set size: 25000, Length: 600), and \textit{FruitFlies} (Training size: 17259, Test set size: 17259, Length: 5000). The FPGA's intermediate computations and inference results are identical to the Aeon baseline. 

We perform two sets of experiments. In the first, the input time series is stored in local memory on both the baseline CPU and FPGA (HBM); in the second, we stream the time series across a network: the FPGA processes the time series directly, while the CPU receives its data via the NIC. 

The CPU baseline experiments were run on a Intel(R) Xeon(R) CPU E5-2640 v3 @ 2.60GHz with 128GiB of DDR4. Latency and throughput was caluclated using overall execution time using Python's built in \textit{time} library and power consumption was estimated using PyRAPL. The FPGA experiments were run in-house on an AMD Xilinx Alveo U280 utilzing the 2022.1 shell, with a Intel(R) Xeon(R) CPU E5-2640 v3 @ 2.60GHz as the host CPU. The test cases and weights are loaded onto off-chip High-Bandwidth Memory (HBM) then loaded onto on-chip BRAM during runtime. Vitis/Vivado 2023.2 was used to develope and verify the design. Resource utilization are given post place and route, and power consumption is measured at runtime by on-board sensors. 

We ran the network experiments on the Open-Cloud Testbed (OCT) \cite{oct, networkfpga} accessed via CloudLabs \cite{CloudLabs}. OCT is a cluster of 24 AMD Xilinx Alveo U280 connected by 100GbE switches, with host Intel(R) Xeon(R) Gold 6226R CPU @ 2.90GHz connected to the same cluster utilizing an Ethernet Controller XL710 for 40GbE. The UltraScale+ Integrated 100G Ethernet Subsystem and AMD's UDP IP \cite{xup_vitis_network_example} provided networking to the FPGA design. The client node utilizes Python Sockets to send UDP packets containing a new time series datapoint as the payload. The inference server (CPU baseline or FPGA) maintains time series window: the arrival of a new datapoint triggers  inference; the server transmits the predicted classes back to the client over UDP. The timing is captured on the client side as the Round-Trip Time (RTT) from when a packet is sent to when the predicted class is received. 

\subsection{Results}

\begin{table}[]
\begin{tabular}{|l|l|l|l|}
\hline
                      & Naive + FP & Naive + bFP & Stream + bFP \\ \hline
Kernel                &           &             &              \\ \hline
Regression Classifier &            &             &              \\ \hline
\end{tabular}
\end{table}




\label{sec:Results}